\documentclass[11pt,a4paper]{article}

% --- Codificación e idioma ---
\usepackage[utf8]{inputenc}
\usepackage[T1]{fontenc}
\usepackage[spanish,es-tabla]{babel}
\usepackage{lmodern}

% --- Geometría y formato general ---
\usepackage[margin=2.5cm]{geometry}
\usepackage{microtype}
\usepackage{parskip}
\usepackage{enumitem}
\setlist{topsep=2pt,itemsep=2pt,parsep=0pt}

% --- Tablas ---
\usepackage{array}
\usepackage{booktabs}
\usepackage{longtable}

% --- Cabeceras y numeración ---
\usepackage{fancyhdr}
\pagestyle{fancy}
\fancyhf{}
\fancyhead[L]{\small Bases de Datos II}
\fancyhead[R]{\small Resumen Teórico --- Triggers y Stored Procedures}
\fancyfoot[C]{\thepage}
\renewcommand{\headrulewidth}{0.4pt}

% --- Color y código ---
\usepackage{xcolor}
\definecolor{codebg}{RGB}{248,248,248}
\definecolor{codeframe}{RGB}{200,200,200}
\definecolor{codekw}{RGB}{0,0,180}
\definecolor{codestr}{RGB}{163,21,21}
\definecolor{codecmt}{RGB}{0,128,0}
\definecolor{notebg}{RGB}{255,250,230}
\definecolor{noteframe}{RGB}{220,180,80}

\usepackage{listings}
\lstdefinelanguage{SQLext}[]{SQL}{
    morekeywords={NUMERIC,VARCHAR,DECIMAL,DATE,TIMESTAMP,TIME,
        BIGINT,VARCHAR2,NUMBER,ENUM,SCHEMA,SEQUENCE,TRIGGER,
        BEFORE,AFTER,EACH,ROW,STATEMENT,NEXTVAL,DUAL,RAISE,EXCEPTION,LOOP,
        DATABASE,DOMAIN,GRANT,REVOKE,ROLE,USER,LOGIN,NOLOGIN,
        OPTION,REFERENCES,RESTRICT,CASCADE,
        AUTO\_INCREMENT,IF,EXISTS,ALTER,SHOW,USE,
        SEARCH\_PATH,EXECUTE,IMMEDIATE,PROCEDURE,FUNCTION,
        DELIMITER,SIGNAL,SQLSTATE,MESSAGE\_TEXT,
        DECLARE,BEGIN,END,RETURNS,LANGUAGE,PLPGSQL,
        OR,REPLACE,RETURN,FOR,IN,OUT,INOUT,
        WITH,FROM,TO,ON,AS,AND,OR,IS,NOT,NULL,
        INSERT,UPDATE,DELETE,SELECT,
        SECURITY,INVOKER,DEFINER,ROUTINE,PROFILE,
        BIGSERIAL,SERIAL,IDENTITY,GENERATED,ALWAYS,
        CONTAINS,READS,SQL,DATA,MODIFIES,COMMENT,
        CONTINUE,EXIT,UNDO,HANDLER,CURSOR,FETCH,OPEN,CLOSE,MOVE,
        SQLEXCEPTION,SQLWARNING,FOUND,NOT,
        REPEAT,UNTIL,WHILE,LEAVE,ITERATE,CASE,WHEN,THEN,ELSE,
        DBMS\_OUTPUT,PUT\_LINE,RAISE\_APPLICATION\_ERROR,SYSDATE,SQLCODE,SQLERRM,
        ROWTYPE,RECORD,SYSTEM,SYSDBA,DEFAULT,TEMPORARY,RULE,INSTEAD,ALSO,NEW,OLD},
    sensitive=false,
    morecomment=[l]{--},
    morestring=[b]'
}
\lstset{
    language=SQLext,
    basicstyle=\ttfamily\footnotesize,
    keywordstyle=\color{codekw}\bfseries,
    stringstyle=\color{codestr},
    commentstyle=\color{codecmt}\itshape,
    backgroundcolor=\color{codebg},
    frame=single,
    rulecolor=\color{codeframe},
    breaklines=true,
    breakatwhitespace=true,
    columns=fullflexible,
    keepspaces=true,
    showstringspaces=false,
    captionpos=b,
    xleftmargin=0.4em,
    xrightmargin=0.4em,
    aboveskip=8pt,
    belowskip=8pt,
    tabsize=2
}

% --- Cajas de nota ---
\usepackage[most]{tcolorbox}
\newtcolorbox{nota}{colback=notebg,colframe=noteframe,arc=2pt,
    left=6pt,right=6pt,top=4pt,bottom=4pt,boxrule=0.6pt}

% --- Hipervínculos ---
\usepackage[hidelinks,bookmarks=true]{hyperref}

\title{Resumen Teórico --- Triggers y Procedimientos Almacenados\\[4pt]
       \large Bases de Datos II --- Práctico N.\textsuperscript{o} 3}
\author{Tomás Rodeghiero}
\date{2026}

\begin{document}

% =====================================================
% Carátula
% =====================================================
\begin{titlepage}
    \centering
    \vspace*{1.5cm}

    {\large \textbf{Universidad Nacional de Río Cuarto}}\\[6pt]
    {\normalsize Facultad de Ciencias Exactas, Físico-Químicas y Naturales}\\[4pt]
    {\normalsize Departamento de Computación}\\[4pt]
    {\normalsize Licenciatura en Ciencias de la Computación}\\[3cm]

    {\Large \textbf{Bases de Datos II}}\\[1.2cm]

    {\Large \textbf{Resumen Teórico}}\\[6pt]
    {\large Triggers y Procedimientos Almacenados\\[2pt]
    (Server Programming) --- Práctico N.\textsuperscript{o} 3}\\[3cm]

    \begin{tabular}{rl}
        \textbf{Alumno:} & Tomás Rodeghiero \\[4pt]
        \textbf{Año:}    & 2026 \\
    \end{tabular}

    \vfill
    {\small Material elaborado a partir del teórico
    \emph{Teorico\_4\_Triggers\_y\_Procedimientos\_Almacenados}.}
\end{titlepage}

\tableofcontents
\newpage

% =====================================================
\section{Introducción}
% =====================================================

Hasta el Práctico 2, todo lo que escribíamos en SQL se ejecutaba como
una sentencia suelta enviada por el cliente. El Práctico 3 introduce
la idea de \textbf{programación del lado servidor} (\emph{server
programming}): bloques de código que viven dentro del motor y se
disparan al ser invocados por un cliente o, automáticamente, ante
ciertos eventos sobre las tablas. Hay dos grandes herramientas:

\begin{description}[leftmargin=2.4cm,style=nextline]
    \item[\textbf{Procedimientos y funciones almacenadas}]
    Programas con SQL embebido, parámetros, variables, control de
    flujo y manejo de errores. Se invocan explícitamente desde el
    cliente con \texttt{CALL} (MySQL/PostgreSQL) o desde un bloque
    \texttt{BEGIN \dots\ END} (Oracle).

    \item[\textbf{Triggers (disparadores)}]
    Procedimientos que el motor ejecuta \emph{automáticamente} cuando
    ocurre un evento sobre una tabla (típicamente
    \texttt{INSERT}/\texttt{UPDATE}/\texttt{DELETE}). Son ideales para
    auditoría, mantener consistencia entre tablas relacionadas y
    normalización de datos antes de su persistencia.
\end{description}

\begin{nota}
\textbf{Idea clave.} El estándar SQL:1999 incorpora \emph{stored
procedures} y \emph{triggers}, pero cada motor mantiene su propio
lenguaje: \textbf{PL/SQL} en Oracle, \textbf{PL/pgSQL} en PostgreSQL,
\textbf{SQL/PSM} (sintaxis SQL:2003) en MySQL. La lógica conceptual
es la misma; lo que cambia es la sintaxis y los detalles de cada
motor.
\end{nota}

% =====================================================
\section{Procedimientos almacenados: visión general}
% =====================================================

Un \textbf{procedimiento almacenado} (\emph{stored procedure}) es un
programa que contiene comandos SQL y se almacena dentro del servidor
de la base de datos. Posee parámetros de entrada y salida y puede ser
invocado por un cliente como si fuera una rutina externa, pero corre
\emph{dentro} del motor.

Los procedimientos están disponibles desde el estándar SQL:1999.
PostgreSQL incluso permite escribirlos en lenguajes externos (C,
Java, Python, etc.). En este resumen nos centramos en los lenguajes
``nativos'' que son los que se usan en el práctico:

\begin{itemize}
    \item \textbf{MySQL} sigue la sintaxis \textbf{SQL:2003} (lenguaje
    \emph{SQL/PSM}).
    \item \textbf{PostgreSQL} usa \textbf{PL/pgSQL}.
    \item \textbf{Oracle} usa \textbf{PL/SQL}.
\end{itemize}

% -----------------------------------------------------
\subsection{Ventajas y desventajas}
% -----------------------------------------------------

\begin{center}
\begin{tabular}{@{}p{3.5cm} p{10.5cm}@{}}
\toprule
\textbf{Ventaja} & \textbf{Por qué importa} \\
\midrule
Rendimiento & Se ejecutan en el motor, sin viaje de datos al
cliente. \\
Compilación previa & Quedan analizados y optimizados al
crearse, no en cada llamada. \\
Menos tráfico de red & Una sola llamada en lugar de N sentencias. \\
Centralización & Forman parte de la base; siempre están
disponibles. \\
Seguridad & Se puede otorgar \texttt{EXECUTE} sobre el
procedimiento sin dar permisos sobre las tablas que toca. \\
Encapsulamiento & El cliente invoca ``\emph{eliminar cliente}''
sin conocer los detalles internos. \\
\bottomrule
\end{tabular}
\end{center}

\paragraph{Desventaja principal.} Cada motor tiene su propio dialecto
y, por lo tanto, migrar lógica entre motores es trabajoso. Es el
precio del poder.

% =====================================================
\section{Procedimientos almacenados en MySQL (SQL/PSM)}
% =====================================================

MySQL incorpora \emph{triggers} y \emph{stored procedures} desde la
versión 5. La sintaxis básica es:

\begin{lstlisting}
CREATE PROCEDURE sp_name ([parameter [, ...]]) [characteristic ...]
   routine_body;

CREATE FUNCTION sp_name ([parameter [, ...]]) RETURNS type
   [characteristic ...] routine_body;

parameter:     [IN | OUT | INOUT] param_name type
characteristic: LANGUAGE SQL
              | { CONTAINS SQL | NO SQL | READS SQL DATA | MODIFIES SQL DATA }
              | SQL SECURITY { DEFINER | INVOKER }
              | COMMENT 'string'
\end{lstlisting}

\paragraph{Parámetros \texttt{IN}/\texttt{OUT}/\texttt{INOUT}.}
\texttt{IN} es entrada (por valor), \texttt{OUT} es de salida
(escrito por la rutina), e \texttt{INOUT} se comporta como ambos.
Si no se especifica el modo, MySQL asume \texttt{IN}.

\paragraph{Características.} \texttt{CONTAINS SQL} indica que el
cuerpo no lee ni modifica datos; \texttt{READS SQL DATA} sólo lee;
\texttt{MODIFIES SQL DATA} puede modificar. \texttt{SQL SECURITY
DEFINER} (default) ejecuta con permisos de quien creó la rutina;
\texttt{INVOKER} con los del invocador.

% -----------------------------------------------------
\subsection{Estructura del lenguaje}
% -----------------------------------------------------

El \emph{routine\_body} se construye con los siguientes elementos:

\begin{itemize}
    \item Sentencias compuestas \texttt{BEGIN \dots\ END} (bloques).
    \item Sentencias \texttt{DECLARE} para variables locales,
    condiciones, \emph{handlers} y cursores.
    \item Asignación de variables (\texttt{SET}, \texttt{SELECT \dots\
    INTO}).
    \item Sentencias de control de flujo: \texttt{IF}, \texttt{CASE},
    \texttt{LOOP}, \texttt{LEAVE}, \texttt{ITERATE}, \texttt{REPEAT},
    \texttt{WHILE}.
    \item Cursores para recorrer resultados de \texttt{SELECT} fila
    por fila.
    \item Manejo de excepciones mediante \emph{handlers}.
\end{itemize}

\paragraph{Bloques \texttt{BEGIN \dots\ END}.} Todo procedimiento no
trivial está delimitado por \texttt{BEGIN \dots\ END}. Pueden
anidarse: cada subbloque tiene su propio scope de variables, sus
declaraciones y sus \emph{handlers}.

% -----------------------------------------------------
\subsection{Variables locales y asignación}
% -----------------------------------------------------

\begin{lstlisting}
-- Declaracion (en el orden correcto: declaraciones siempre arriba del bloque).
DECLARE var_name type [DEFAULT value];
DECLARE x INT;
DECLARE saldo DECIMAL(12,2) DEFAULT 0;

-- Asignacion clasica.
SET x = 10;

-- Asignacion via SELECT.
SELECT MAX(precio) INTO x FROM producto;
\end{lstlisting}

% -----------------------------------------------------
\subsection{\emph{Handlers} y manejo de excepciones}
% -----------------------------------------------------

Los \emph{handlers} reaccionan ante condiciones (errores o avisos)
durante la ejecución:

\begin{lstlisting}
DECLARE handler_type HANDLER FOR condition_value [, ...] sp_statement;

handler_type:    CONTINUE | EXIT | UNDO
condition_value: SQLSTATE sqlstate_value
              | condition_name
              | SQLWARNING
              | NOT FOUND
              | SQLEXCEPTION
              | mysql_error_code
\end{lstlisting}

\paragraph{Tipos de \emph{handler}.}

\begin{itemize}
    \item \texttt{CONTINUE}: ejecuta la sentencia del handler y
    continúa con la siguiente sentencia del programa.
    \item \texttt{EXIT}: ejecuta el handler y \emph{termina} el
    bloque que lo contiene.
    \item \texttt{UNDO}: deshace lo hecho dentro del bloque (no
    siempre soportado).
\end{itemize}

\paragraph{Tipos de condición.}

\begin{itemize}
    \item \texttt{SQLSTATE 'xxxxx'}: estándar, comparten todos los
    motores.
    \item \texttt{SQLWARNING}: cualquier \texttt{SQLSTATE} que empiece
    con \texttt{01}.
    \item \texttt{NOT FOUND}: cualquier \texttt{SQLSTATE} que empiece
    con \texttt{02}. Es el típico ``ya no hay más filas'' de un
    cursor.
    \item \texttt{SQLEXCEPTION}: todos los \texttt{SQLSTATE}
    restantes.
    \item \texttt{mysql\_error\_code}: códigos específicos de MySQL.
\end{itemize}

\paragraph{Lanzar excepciones del usuario.} Desde MySQL 5.5:

\begin{lstlisting}
SIGNAL SQLSTATE '45000'
    SET MESSAGE_TEXT = 'lanzo una excepcion';
\end{lstlisting}

% -----------------------------------------------------
\subsection{Cursores}
% -----------------------------------------------------

Un \emph{cursor} permite recorrer fila por fila el resultado de una
consulta. Su uso típico tiene cuatro etapas: \texttt{DECLARE},
\texttt{OPEN}, repetidos \texttt{FETCH}, \texttt{CLOSE}.

\begin{lstlisting}
-- Declaracion.
DECLARE cursor_name CURSOR FOR select_statement;

-- Operaciones.
OPEN  cursor_name;
FETCH cursor_name INTO var1 [, var2, ...];
CLOSE cursor_name;
\end{lstlisting}

% -----------------------------------------------------
\subsection{Sentencias de control de flujo}
% -----------------------------------------------------

\begin{itemize}
    \item \texttt{IF cond THEN \dots\ ELSE \dots\ END IF;}
    \item \texttt{CASE \dots\ WHEN \dots\ THEN \dots\ END CASE;}
    \item \texttt{LOOP \dots\ END LOOP;} (con \texttt{LEAVE} para
    salir e \texttt{ITERATE} para volver al principio).
    \item \texttt{REPEAT \dots\ UNTIL cond END REPEAT;}
    \item \texttt{WHILE cond DO \dots\ END WHILE;}
\end{itemize}

\paragraph{REPEAT (semántica).} Las sentencias del bloque se ejecutan
\textbf{al menos una vez}, y se repiten hasta que la condición
\texttt{search\_condition} se vuelva verdadera.

\begin{lstlisting}
REPEAT
   statement_list
UNTIL search_condition
END REPEAT;
\end{lstlisting}

% -----------------------------------------------------
\subsection{Ejemplo completo: cursor + handler + REPEAT}
% -----------------------------------------------------

Recorre \texttt{tabla1} filtrando por un parámetro y, según una
comparación entre dos columnas, inserta en \texttt{tabla2} o en
\texttt{tabla3}:

\begin{lstlisting}
CREATE PROCEDURE cursordemo(IN parametro INT)
BEGIN
    DECLARE done INT DEFAULT 0;
    DECLARE a1 CHAR(16);
    DECLARE b1, c1 INT;

    DECLARE cursor1 CURSOR FOR
        SELECT a, b, c FROM tabla1 WHERE id < parametro;

    -- Cuando NOT FOUND (SQLSTATE 02000), termina el recorrido.
    DECLARE CONTINUE HANDLER FOR SQLSTATE '02000'
    BEGIN
        SET done = 1;
    END;

    OPEN cursor1;
    REPEAT
        FETCH cursor1 INTO a1, b1, c1;
        IF NOT done THEN
            IF b1 < c1 THEN
                INSERT INTO tabla2 VALUES (a1, b1);
            ELSE
                INSERT INTO tabla3 VALUES (a1, c1);
            END IF;
        END IF;
    UNTIL done END REPEAT;
    CLOSE cursor1;
END;
\end{lstlisting}

\paragraph{Lectura del ejemplo.} El \emph{handler} de tipo
\texttt{CONTINUE} captura el aviso ``no hay más filas'' y pone
\texttt{done = 1}. El bucle \texttt{REPEAT} hace \texttt{FETCH}
dentro y termina cuando \texttt{done} pasa a 1.

% -----------------------------------------------------
\subsection{Modificación e invocación}
% -----------------------------------------------------

Para cambiar las características de una rutina existente (por ejemplo,
cambiar de \texttt{DEFINER} a \texttt{INVOKER}, o agregar
\texttt{COMMENT}):

\begin{lstlisting}
ALTER {PROCEDURE | FUNCTION} sp_name [characteristic ...];
\end{lstlisting}

Para invocar un procedimiento se usa \texttt{CALL}:

\begin{lstlisting}
CALL nombreProcedimiento([parametroActual [, ...]]);
CALL cursordemo(2);
\end{lstlisting}

% =====================================================
\section{Procedimientos y funciones en PostgreSQL (PL/pgSQL)}
% =====================================================

PostgreSQL implementa funciones (\texttt{CREATE FUNCTION}) y, desde la
versión 11, también \emph{procedures} (\texttt{CREATE PROCEDURE}).
Pueden escribirse en distintos lenguajes; el de uso más común es
\textbf{PL/pgSQL}, orientado a bloques.

% -----------------------------------------------------
\subsection{Estructura de PL/pgSQL}
% -----------------------------------------------------

\begin{lstlisting}
CREATE [OR REPLACE] FUNCTION nombreFuncion ([parametro [, ...]])
RETURNS tipoRetorno AS $$
[DECLARE
    declaracion_de_variables_locales]
BEGIN
    sentencias_propias_de_la_funcion;
[EXCEPTION
    WHEN condicion [OR condicion ...]
        THEN sentencias_para_manejar_la_excepcion
    [WHEN condicion [OR condicion ...]
        THEN sentencias_para_manejar_la_excepcion ...]
]
END;
$$ LANGUAGE 'plpgsql';
\end{lstlisting}

\paragraph{Delimitadores \$\$.} El cuerpo entero se encierra entre
delimitadores ``dólar comillado''. PostgreSQL acepta
\texttt{\$nombre\$ \dots\ \$nombre\$} si se necesita anidar.

% -----------------------------------------------------
\subsection{Parámetros y variables}
% -----------------------------------------------------

\begin{lstlisting}
-- Parametros: [IN | OUT | INOUT] [nombre] tipo
CREATE FUNCTION sum_n_product(int, y int, OUT sum int, OUT prod int) AS $$
BEGIN
    sum  := $1 + y;     -- $1 referencia al primer parametro sin nombre.
    prod := $1 * y;
END;
$$ LANGUAGE plpgsql;
\end{lstlisting}

Si una función se declara con uno o más \texttt{OUT}, PostgreSQL la
trata como devolviendo una fila con esas columnas. Se invoca con un
\texttt{SELECT}.

\paragraph{Declaración de variables.}

\begin{lstlisting}
nombre [CONSTANT] tipo [NOT NULL] [{ DEFAULT | := } expresion];

-- Alias de un parametro posicional.
nombreAlias ALIAS FOR $n;

-- Asignacion.
nombreVariable := expresion;
\end{lstlisting}

% -----------------------------------------------------
\subsection{Bloques anidados y control de flujo}
% -----------------------------------------------------

Igual que en MySQL, los bloques se delimitan con
\texttt{BEGIN \dots\ END} y pueden anidarse, cada uno con sus
propios manejadores de excepciones.

PL/pgSQL provee:

\begin{itemize}
    \item \texttt{LOOP \dots\ END LOOP;}
    \item \texttt{CASE}
    \item \texttt{EXIT} (sale del bloque/ciclo)
    \item \texttt{IF \dots\ THEN \dots\ ELSE \dots\ END IF;}
    \item \texttt{FOR}
    \item \texttt{WHILE}
\end{itemize}

\paragraph{Sentencia \texttt{FOR} numérica.}

\begin{lstlisting}
FOR variable IN inicio..fin LOOP
    cuerpo_del_for;
END LOOP;

-- Iteracion en sentido inverso.
FOR variable IN REVERSE fin..inicio LOOP
    cuerpo_del_for;
END LOOP;

-- Ejemplo: factorial.
FOR inc IN 1..10 LOOP
    factorial := factorial * inc;
END LOOP;
\end{lstlisting}

\paragraph{Sentencia \texttt{FOR} sobre filas (cursor implícito).}
PL/pgSQL permite recorrer un \texttt{SELECT} sin abrir el cursor
manualmente:

\begin{lstlisting}
DECLARE
    registro RECORD;          -- registro generico
BEGIN
    FOR registro IN SELECT * FROM productos LOOP
        stock := registro.stock_actual + 100;
    END LOOP;
END;
\end{lstlisting}

\begin{nota}
\textbf{Tipos de registro.} Se puede declarar un registro
\texttt{RECORD} (genérico) o ligado al esquema de una tabla con
\texttt{persona\%ROWTYPE}. Esto último brinda chequeo de tipos por
columna.
\end{nota}

% -----------------------------------------------------
\subsection{Cursores explícitos}
% -----------------------------------------------------

\begin{lstlisting}
DECLARE nombre CURSOR FOR sentenciaSelect;

-- Operaciones: FETCH, CLOSE, MOVE.
-- La variable booleana FOUND indica si la ultima operacion devolvio fila.
\end{lstlisting}

% -----------------------------------------------------
\subsection{Excepciones}
% -----------------------------------------------------

PostgreSQL permite lanzar excepciones del usuario:

\begin{lstlisting}
RAISE EXCEPTION 'el mensaje %', valor;
\end{lstlisting}

Y capturarlas con bloques \texttt{EXCEPTION \dots\ WHEN}:

\begin{lstlisting}
BEGIN
    -- ...
EXCEPTION
    WHEN condicion [OR condicion ...] THEN
        sentencias_para_manejar_la_excepcion;
END;
\end{lstlisting}

% -----------------------------------------------------
\subsection{Invocación}
% -----------------------------------------------------

A diferencia de MySQL, en PostgreSQL las funciones se invocan con
\texttt{SELECT}:

\begin{lstlisting}
SELECT nombreProcedimiento([parametroActual [, ...]]);
SELECT procedimiento(2);
\end{lstlisting}

Si la función tiene parámetros \texttt{OUT}, conviene
\texttt{SELECT * FROM funcion(\dots)} para ver todas las columnas
devueltas.

% =====================================================
\section{Procedimientos y funciones en Oracle (PL/SQL)}
% =====================================================

Oracle ofrece un lenguaje muy potente, \textbf{PL/SQL}, con
estructura semejante a la de PL/pgSQL.

% -----------------------------------------------------
\subsection{Sintaxis general}
% -----------------------------------------------------

\paragraph{Procedimiento.}

\begin{lstlisting}
CREATE [OR REPLACE] PROCEDURE nombreProc ([parametro [, ...]])
AS
    [DECLARE]
        declaraciones
BEGIN
    programa;
[EXCEPTION
    tratamiento_de_excepciones]
END;
/
\end{lstlisting}

\paragraph{Función.}

\begin{lstlisting}
CREATE [OR REPLACE] FUNCTION nombreFunc ([parametro [, ...]])
RETURN tipo
AS
    [declaraciones]
BEGIN
    programa;
[EXCEPTION
    tratamiento_de_excepciones]
END;
/
\end{lstlisting}

\paragraph{Parámetros.}

\begin{lstlisting}
nombreParametro [IN | OUT | IN OUT] tipo [DEFAULT expr]
\end{lstlisting}

\paragraph{Variables.}

\begin{lstlisting}
nombreVariable tipo [NOT NULL] [:= valor | DEFAULT valor];
nombreVariable := expresion;       -- asignacion
\end{lstlisting}

% -----------------------------------------------------
\subsection{Excepciones en Oracle}
% -----------------------------------------------------

Oracle permite lanzar excepciones del usuario con
\texttt{RAISE\_APPLICATION\_ERROR}, cuyo número debe estar en el
rango $[-20999, -20000]$:

\begin{lstlisting}
RAISE_APPLICATION_ERROR(Error_number INT, texto TEXT);

-- Ejemplo.
RAISE_APPLICATION_ERROR(-20011, 'prueba de excepcion');
\end{lstlisting}

También se puede declarar una excepción con nombre y lanzarla con
\texttt{RAISE}.

\paragraph{Captura de excepciones.} Con \texttt{WHEN OTHERS THEN \dots},
y las funciones \texttt{SQLCODE} y \texttt{SQLERRM} para obtener el
código y el mensaje:

\begin{lstlisting}
DECLARE
    err_num NUMBER;
    err_msg VARCHAR2(255);
    result  NUMBER;
BEGIN
    RAISE_APPLICATION_ERROR(-20011, 'prueba de excepcion');
EXCEPTION
    WHEN OTHERS THEN
        err_num := SQLCODE;
        err_msg := SQLERRM;
        DBMS_OUTPUT.PUT_LINE('Error: ' || TO_CHAR(err_num));
        DBMS_OUTPUT.PUT_LINE(err_msg);
END;
\end{lstlisting}

% -----------------------------------------------------
\subsection{Cursores en PL/SQL}
% -----------------------------------------------------

La sintaxis de \texttt{DECLARE / OPEN / FETCH / CLOSE} es la misma
que en MySQL. La diferencia interesante son los \emph{atributos} que
expone cada cursor:

\begin{itemize}
    \item \texttt{\%FOUND}: \texttt{true} si el último \texttt{FETCH}
    devolvió fila.
    \item \texttt{\%NOTFOUND}: complemento del anterior.
    \item \texttt{\%ISOPEN}: \texttt{true} si el cursor está abierto.
    \item \texttt{\%ROWCOUNT}: cantidad de filas leídas hasta el
    momento.
\end{itemize}

\paragraph{Ejemplo de manejo de cursor.}

\begin{lstlisting}
CREATE OR REPLACE PROCEDURE listar IS
    CURSOR cproductos IS SELECT nProd, precio FROM productos;
    regProducto cproductos%ROWTYPE;
BEGIN
    OPEN cproductos;
    FETCH cproductos INTO regProducto;
    WHILE cproductos%FOUND LOOP
        DBMS_OUTPUT.PUT_LINE(regProducto.precio);
        FETCH cproductos INTO regProducto;
    END LOOP;
    CLOSE cproductos;
END;
/
\end{lstlisting}

% -----------------------------------------------------
\subsection{Invocación en Oracle}
% -----------------------------------------------------

Se invoca dentro de un bloque anónimo:

\begin{lstlisting}
[DECLARE
    declaraciones]
BEGIN
    nombreProc([parametroActual [, ...]]);
END;
/

-- Ejemplo.
DECLARE
    nombre_parametro INT;
BEGIN
    nombre_parametro := 3;
    pruebaProc(3, nombre_parametro);
END;
/
\end{lstlisting}

% =====================================================
\section{Triggers (disparadores)}
% =====================================================

Un \textbf{trigger} es un procedimiento que el motor ejecuta
\emph{automáticamente} cuando se cumple una \emph{condición de
disparo}: una operación DML (\texttt{INSERT}, \texttt{UPDATE} o
\texttt{DELETE}) sobre una tabla específica. Se utilizan
principalmente para:

\begin{itemize}
    \item generar información de auditoría;
    \item mantener consistencia entre tablas;
    \item normalizar datos antes de su persistencia (mayúsculas,
    \texttt{TRIM}, etc.);
    \item completar campos derivados (autonumerado en motores que no
    lo traen nativo).
\end{itemize}

Para diseñar un trigger hay que definir:

\begin{enumerate}
    \item la \textbf{condición de disparo} (qué evento, sobre qué
    tabla, antes o después);
    \item las \textbf{acciones} a ejecutar cuando se dispara.
\end{enumerate}

% -----------------------------------------------------
\subsection{Variables contextuales \texttt{NEW} y \texttt{OLD}}
% -----------------------------------------------------

Dentro del cuerpo del trigger, dos variables especiales dan acceso a
la fila afectada:

\begin{itemize}
    \item \texttt{NEW}: contiene los \emph{nuevos} valores. Disponible
    en \texttt{INSERT} y \texttt{UPDATE}.
    \item \texttt{OLD}: contiene los \emph{viejos} valores. Disponible
    en \texttt{UPDATE} y \texttt{DELETE}.
\end{itemize}

En PostgreSQL y MySQL se referencian directamente como
\texttt{NEW.col}/\texttt{OLD.col}; en Oracle el cuerpo PL/SQL las
referencia con dos puntos: \texttt{:NEW.col}/\texttt{:OLD.col}.

% -----------------------------------------------------
\subsection{Estándar y dialectos por motor}
% -----------------------------------------------------

\begin{itemize}
    \item SQL-92 \emph{no} incluye triggers; cada motor tenía su
    propia forma.
    \item SQL:1999 incorporó la sintaxis estándar, pero los motores
    siguen ofreciendo extensiones.
    \item El \emph{lenguaje} del cuerpo del trigger es el mismo que
    para procedimientos almacenados (SQL/PSM, PL/pgSQL o PL/SQL).
\end{itemize}

% -----------------------------------------------------
\subsection{Triggers en MySQL}
% -----------------------------------------------------

Disponibles desde MySQL 5.0.2.

\begin{lstlisting}
CREATE TRIGGER nombre_trigger momento_disparo evento_disparo
    ON nombre_tabla
    FOR EACH ROW
    sentencia_trigger;

-- momento_disparo: BEFORE | AFTER
-- evento_disparo:  INSERT | UPDATE | DELETE
\end{lstlisting}

\paragraph{Ejemplo: auditar bajas en \texttt{articulos}.}

\begin{lstlisting}
delimiter $$
CREATE TRIGGER trigger_baja_articulos
AFTER DELETE ON articulos
FOR EACH ROW
BEGIN
    INSERT INTO articulos_bajas
    VALUES (OLD.nart, OLD.descr, OLD.precio, CURRENT_USER(), NOW());
END;
$$
delimiter ;
\end{lstlisting}

\begin{nota}
\textbf{\texttt{CURRENT\_USER()} vs usuario de aplicación.} En MySQL,
\texttt{CURRENT\_USER()} devuelve el usuario con el que se estableció
la conexión (típicamente uno técnico de la aplicación, como
\texttt{webapp@\%}). Casi siempre la auditoría busca otra cosa: el
usuario \emph{de la aplicación cliente} (la persona logueada). Eso
debe registrarse en una variable de sesión:
\end{nota}

\begin{lstlisting}
-- Setear desde el cliente al iniciar la sesion.
SET @id_usuario = 4;

-- El trigger lee la variable de sesion.
INSERT INTO AUDITORIA VALUES (0, @id_usuario, NEW.fecha);
\end{lstlisting}

% -----------------------------------------------------
\subsection{Triggers en PostgreSQL}
% -----------------------------------------------------

PostgreSQL desacopla \emph{trigger} y \emph{función}: primero se crea
una función que devuelve \texttt{TRIGGER}, después se asocia esa
función a una tabla con \texttt{CREATE TRIGGER}.

\begin{lstlisting}
CREATE TRIGGER nombre {BEFORE | AFTER}
    {evento [OR ...]}
    ON tabla [FOR [EACH] {ROW | STATEMENT}]
    EXECUTE PROCEDURE nombreFunc(argumentos);
\end{lstlisting}

\paragraph{Ejemplo de auditoría.}

\begin{lstlisting}
-- 1) Funcion que se ejecutara al disparar el trigger.
CREATE FUNCTION funcion_auditoria() RETURNS TRIGGER AS $$
BEGIN
    INSERT INTO auditoria
    VALUES (1, OLD.nart, NOW(), OLD.cant - NEW.cant);
    RETURN NEW;
END;
$$ LANGUAGE 'plpgsql';

-- 2) Asociar el trigger a la tabla.
CREATE TRIGGER trigger_auditoria
AFTER UPDATE ON articulos
FOR EACH ROW
EXECUTE PROCEDURE funcion_auditoria();
\end{lstlisting}

\paragraph{Diferencia \texttt{ROW} vs \texttt{STATEMENT}.} Con
\texttt{FOR EACH ROW} la función se ejecuta una vez por cada fila
afectada; con \texttt{FOR EACH STATEMENT}, una sola vez por sentencia.

% -----------------------------------------------------
\subsection{Triggers en Oracle}
% -----------------------------------------------------

\begin{lstlisting}
CREATE [OR REPLACE] TRIGGER [esquema.]nombre
{BEFORE | AFTER | INSTEAD OF}
{ dml_event_clause
| ddl_event [OR ddl_eventN]...
| database_event [OR database_eventN]...}
ON {tabla | view | [esquema.]SCHEMA | DATABASE}
[FOR EACH ROW]
[WHEN (condicion)]
{pl/sql_block | call_procedure_statement};
\end{lstlisting}

\paragraph{Ejemplo (equivalente al de MySQL).}

\begin{lstlisting}
CREATE TRIGGER trigger_baja_articulos
AFTER DELETE ON articulos
FOR EACH ROW
BEGIN
    INSERT INTO articulos_bajas
    VALUES (:OLD.nart, :OLD.descr, :OLD.precio, USER, SYSDATE);
END;
/
\end{lstlisting}

\paragraph{\texttt{INSTEAD OF}.} Es propio de Oracle/PostgreSQL y se
asocia a \emph{vistas}: en lugar de ejecutar la operación sobre la
vista (que muchas veces no es actualizable), el trigger redirige la
acción a las tablas base.

\paragraph{\texttt{WHEN (condicion)}.} Permite filtrar el disparo:
sólo se ejecuta el cuerpo si la condición sobre \texttt{NEW}/\texttt{OLD}
es verdadera. Es lo que permite, por ejemplo, autonumerar sólo cuando
el cliente no provee un valor explícito:

\begin{lstlisting}
CREATE OR REPLACE TRIGGER bi_parquimetro_autonumerado
BEFORE INSERT ON parquimetro
FOR EACH ROW
WHEN (NEW.id_parquimetro IS NULL)
BEGIN
    SELECT seq_parquimetro.NEXTVAL
      INTO :NEW.id_parquimetro
      FROM dual;
END;
/
\end{lstlisting}

% =====================================================
\section{Reglas de reescritura en PostgreSQL}
% =====================================================

PostgreSQL incluye un mecanismo extra, distinto de los triggers,
llamado \textbf{sistema de reglas de reescritura}
(\texttt{CREATE RULE}). En lugar de ejecutar código adicional cuando
ocurre un evento, \emph{modifica el query} antes de ejecutarlo.

\paragraph{Las vistas son reglas.} En PostgreSQL,
\texttt{CREATE VIEW vista AS SELECT * FROM tabla;} es equivalente a:

\begin{lstlisting}
CREATE TABLE vista (lista_de_columnas);
CREATE RULE ReescribeVista AS ON SELECT TO vista
    DO INSTEAD SELECT * FROM tabla;
\end{lstlisting}

es decir, una tabla vacía más una regla que reescribe los
\texttt{SELECT} a esa tabla.

\paragraph{Reglas de actualización.} Permiten interceptar
\texttt{INSERT}, \texttt{UPDATE} o \texttt{DELETE}:

\begin{lstlisting}
CREATE [OR REPLACE] RULE nombre
    AS ON evento TO tabla
    [WHERE condicion]
    DO [ALSO | INSTEAD]
    {NOTHING | comando | (comando; comando ...)};
\end{lstlisting}

\paragraph{Ejemplo.} Auditoría de cambios de precio:

\begin{lstlisting}
CREATE RULE log_precio_productos
AS ON UPDATE TO productos
WHERE NEW.precio <> OLD.precio
DO ALSO
INSERT INTO productos_log
VALUES (NEW.codigo, NEW.nombre, OLD.precio, NEW.precio,
        current_user, current_timestamp);
\end{lstlisting}

% =====================================================
\section{Triggers vs procedimientos}
% =====================================================

\begin{center}
\begin{tabular}{@{}p{4cm} p{4.5cm} p{5.5cm}@{}}
\toprule
& \textbf{Procedimiento almacenado} & \textbf{Trigger} \\
\midrule
Quién lo dispara &
El cliente, explícitamente, con \texttt{CALL}/\texttt{SELECT} &
El motor, \emph{automáticamente}, ante un evento DML \\
\addlinespace
Parámetros & Sí (\texttt{IN}, \texttt{OUT}, \texttt{INOUT}) &
No: el ``input'' es la fila afectada (\texttt{NEW}/\texttt{OLD}) \\
\addlinespace
Frecuencia & Una vez por invocación &
Por fila o por sentencia, según la configuración \\
\addlinespace
Uso típico & Lógica de negocio invocable, encapsulamiento &
Auditoría, integridad cruzada, normalización, autonumerado \\
\addlinespace
Aborta operación & Sólo levantando excepción &
Sí, levantando excepción en \texttt{BEFORE} \\
\bottomrule
\end{tabular}
\end{center}

% =====================================================
\section{Conexión con el Práctico 3}
% =====================================================

\begin{itemize}
    \item \textbf{Ejercicio 1 (MySQL/PostgreSQL)}: agrega
    \texttt{cantidad\_maxima\_items} a \texttt{factura} y la valida
    con triggers \texttt{BEFORE INSERT}/\texttt{UPDATE}, levantando
    excepción con \texttt{SIGNAL SQLSTATE '45000'} (MySQL) o
    \texttt{RAISE EXCEPTION} (PostgreSQL) si se supera el máximo.
    \item \textbf{Ejercicio 2 (PostgreSQL)}: función con parámetros
    \texttt{IN}/\texttt{OUT} que recorre las cotizaciones con un
    \texttt{FOR \dots\ LOOP} (sin usar \texttt{MAX}). Aplica
    directamente la \emph{sentencia FOR sobre filas} del teórico.
    \item \textbf{Ejercicio 3 (Oracle)}: dos procedimientos
    almacenados con \texttt{IN}/\texttt{OUT}, validación con
    \texttt{RAISE\_APPLICATION\_ERROR}, manejo de
    \texttt{NO\_DATA\_FOUND} y salida con
    \texttt{DBMS\_OUTPUT.PUT\_LINE}.
    \item \textbf{Ejercicio 4 (MySQL)}: procedimiento creado por
    \texttt{vendedor2} con \texttt{SQL SECURITY INVOKER}; conecta
    DCL del Práctico 2 con SQL/PSM del Práctico 3.
    \item \textbf{Ejercicio 5}: triggers de tres tipos:
    \emph{normalización} (\texttt{BEFORE INSERT}, \texttt{UPPER}),
    \emph{auditoría} (\texttt{AFTER UPDATE}, leyendo el usuario de
    aplicación de una variable de sesión) y \emph{autonumerado}
    (\texttt{AUTO\_INCREMENT} en MySQL, \texttt{SEQUENCE + TRIGGER
    BEFORE INSERT WHEN (NEW.id IS NULL)} en Oracle).
\end{itemize}

% =====================================================
\section{Ideas clave para repasar antes del parcial}
% =====================================================

\begin{enumerate}
    \item \textbf{Procedimiento vs trigger.} El primero se invoca; el
    segundo se dispara automáticamente.
    \item \textbf{Cada motor, su lenguaje}: SQL/PSM (MySQL), PL/pgSQL
    (PostgreSQL), PL/SQL (Oracle).
    \item \textbf{Bloque \texttt{BEGIN \dots\ END}} es la unidad básica
    de ejecución y de scope; puede anidarse.
    \item \textbf{\texttt{DECLARE} primero, código después}: las
    variables, condiciones, \emph{handlers} y cursores se declaran al
    inicio de cada bloque.
    \item \textbf{Parámetros}: \texttt{IN} (entrada), \texttt{OUT}
    (salida), \texttt{INOUT} (ambos). En PostgreSQL, una función con
    \texttt{OUT} se invoca con \texttt{SELECT * FROM funcion(\dots)}.
    \item \textbf{Excepciones} (cómo lanzarlas):
    \texttt{SIGNAL SQLSTATE '45000'} (MySQL), \texttt{RAISE EXCEPTION
    '\dots'} (PostgreSQL), \texttt{RAISE\_APPLICATION\_ERROR(-20\dots,
    '\dots')} (Oracle).
    \item \textbf{Excepciones} (cómo capturarlas): \emph{handlers}
    \texttt{CONTINUE}/\texttt{EXIT}/\texttt{UNDO} en MySQL, bloques
    \texttt{EXCEPTION WHEN \dots\ THEN} en PostgreSQL/Oracle.
    \item \textbf{Cursores}: ciclo \texttt{DECLARE} $\to$
    \texttt{OPEN} $\to$ \texttt{FETCH (en bucle)} $\to$
    \texttt{CLOSE}. La condición \texttt{NOT FOUND} (\texttt{SQLSTATE
    '02000'}) o el atributo \texttt{\%NOTFOUND} indican fin de
    iteración.
    \item \textbf{Cursor implícito en PL/pgSQL}: el bucle
    \texttt{FOR registro IN SELECT \dots\ LOOP \dots\ END LOOP}
    ahorra el \texttt{OPEN}/\texttt{FETCH}/\texttt{CLOSE} explícitos.
    \item \textbf{Triggers} se diseñan con dos elementos: condición
    de disparo (BEFORE/AFTER + INSERT/UPDATE/DELETE) y cuerpo.
    Dentro: \texttt{NEW}/\texttt{OLD} (con \texttt{:} en Oracle).
    \item \textbf{Auditar el usuario de aplicación} requiere
    variables de sesión, no \texttt{CURRENT\_USER()}.
    \item \textbf{PostgreSQL separa función y trigger}: la función
    \texttt{RETURNS TRIGGER}; el trigger se asocia con
    \texttt{CREATE TRIGGER \dots\ EXECUTE PROCEDURE func()}.
    \item \textbf{\texttt{INSTEAD OF}} es propio de vistas (Oracle y
    PostgreSQL): redirige una operación que se ejecuta sobre una
    vista hacia las tablas base.
    \item \textbf{Reglas de reescritura} son una herramienta extra de
    PostgreSQL: \emph{modifican el query}, no agregan código de
    ejecución.
    \item \textbf{\texttt{ALTER PROCEDURE/FUNCTION}} cambia
    \emph{características} (\texttt{SQL SECURITY}, \texttt{COMMENT},
    etc.); para cambiar el cuerpo se usa \texttt{CREATE OR REPLACE}.
\end{enumerate}

% =====================================================
\section*{Resumen final}
% =====================================================
\addcontentsline{toc}{section}{Resumen final}

Los procedimientos almacenados y los triggers son las herramientas con
las que la lógica de negocio se acomoda dentro del motor de base de
datos. Compartiendo el mismo lenguaje del lado servidor en cada motor
--- SQL/PSM, PL/pgSQL o PL/SQL --- permiten escribir código con
parámetros, variables, control de flujo, cursores y excepciones,
ejecutado lo más cerca posible de los datos.

Si quedan firmes los cinco bloques --- \emph{procedimiento vs
función}, \emph{parámetros \texttt{IN}/\texttt{OUT}}, \emph{control
de flujo y cursores}, \emph{manejo de excepciones} y \emph{triggers
con \texttt{NEW}/\texttt{OLD}} ---, el Práctico 3 se traduce
directamente: cada inciso es una combinación de uno o dos de esos
bloques sobre las bases de datos heredadas de los prácticos
anteriores. La gracia no está sólo en escribir el código, sino en
elegir bien cuándo usar trigger, cuándo procedimiento y cuándo
ninguna de las dos (porque alcanza con un \texttt{CHECK} o una
\texttt{FOREIGN KEY}).

\end{document}
